\section*{Bab \ref{ch:g8:negprod} --- Mengalikan Bilangan Negatif}

\begin{solution}{exo:g8:negprod:1}
$(-8) \times (+7) = -56$; \quad
$(-9) \times (-6) = +54$; \quad
$(+12) \times (-0.5) = -6$; \quad
$(-1) \times (-1) \times (-1) = -1$ (tiga faktor negatif: \omterm{def:g3:numbers:evenodd}{ganjil}).
\end{solution}

\begin{solution}{exo:g8:negprod:2}
$(-63) \div (+9) = -7$; \quad
$(-8) \div (-16) = +0.5$; \quad
$\dfrac{45}{-5} = -9$; \quad
$\dfrac{-7.2}{-8} = +0.9$.
\end{solution}

\begin{solution}{exo:g8:negprod:3}
Tiga faktor negatif: \emph{negatif}. \quad
$(-1)^{10}$: sepuluh faktor negatif, \omterm{def:g3:numbers:evenodd}{genap}: \emph{positif} (nilainya
$1$). \quad
$(-2)^7$: tujuh faktor, \omterm{def:g3:numbers:evenodd}{ganjil}: \emph{negatif}. \quad
$(-4) \times 0 \times (-6) = 0$: bukan positif, bukan pula negatif.
\end{solution}

\begin{solution}{exo:g8:negprod:4}
$(-2)^4 = +16$; \quad $-2^4 = -16$ (minusnya tidak dikuadratkan); \quad
$(-3)^3 = -27$; \quad $(-1)^{2026} = +1$ (pangkatnya \omterm{def:g3:numbers:evenodd}{genap}).
\end{solution}

\begin{solution}{exo:g8:negprod:5}
$7 + 2 \times (-6) = 7 + (-12) = -5$.

$(-4) \times (5 - 9) = (-4) \times (-4) = +16$.

$18 \div (-3) - (-2) = -6 + 2 = -4$.
\end{solution}

\begin{solution}{exo:g8:negprod:6}
$\dfrac{(-5) \times (+8)}{-4} = \dfrac{-40}{-4} = +10$.

$\dfrac{(-3) \times (-6)}{(-2) \times (+9)} = \dfrac{+18}{-18} = -1$.
\end{solution}

\begin{solution}{exo:g8:negprod:7}
Untuk $x = -3$: \ $5x = -15$; \quad $x^2 = (-3)^2 = 9$; \quad
$-x^2 = -9$; \quad $2x^2 + 4x = 18 - 12 = 6$; \quad
$(2x)^2 = (-6)^2 = 36$.
\end{solution}

\begin{solution}{exo:g8:negprod:8}
$(-6) \times (-7) = 42$; \quad
$10 \div (-4) = -2.5$; \quad
$(-5) \times \frac15 = -1$, jadi bilangan yang hilang adalah
$+\frac15 = 0.2$.
\end{solution}

\begin{solution}{exo:g8:negprod:9}
\emph{1. Salah}: $0.5 \times 0.5 = 0.25$ lebih kecil daripada
$0.5 + 0.5 = 1$ (atau: $2 \times 3 = 6 > 5$, tetapi pernyataannya
harus selalu berlaku).

\emph{2. Benar}: kuadrat adalah \omterm{def:g2:mult:def}{hasil kali} dua bilangan yang tandanya
sama, jadi positif atau \omterm{def:g1:counting:zero}{nol}.

\emph{3. Salah}: $(-2) \times (-3) \times 4 = +24$ bernilai positif
dengan dua faktor negatif.
\end{solution}

\begin{solution}{exo:g8:negprod:10}
Misalkan $r$ banyaknya jawaban yang benar; maka $20 - r$ jawaban salah, dan
\[
5r - 3(20 - r) = 36
\ \Longrightarrow\
5r - 60 + 3r = 36
\ \Longrightarrow\
8r = 96
\ \Longrightarrow\
r = 12 .
\]
Dua belas jawaban yang benar (dan delapan yang salah: periksa,
$60 - 24 = 36$).
\end{solution}

\begin{solution}{exo:g8:negprod:11}
\omterm{def:g1:counting:zero}{Nol} dikalikan apa pun adalah \omterm{def:g1:counting:zero}{nol}, jadi $(-1) \times 0 = 0$. Dengan
menulis $0 = 1 + (-1)$ lalu menyebarkannya:
\[
0 = (-1) \times \bigl(1 + (-1)\bigr)
= (-1) \times 1 + (-1) \times (-1)
= -1 + (-1) \times (-1).
\]
Bilangan $(-1) \times (-1)$ yang ditambahkan ke $-1$ memberi $0$: ia
pastilah \omterm{def:g7:negatives:def}{lawan} dari $-1$, yaitu $+1$. Sifat distributif tidak
meninggalkan pilihan lain.
\end{solution}

\begin{solution}{pb:g8:negprod:1}
\textbf{1.} Karena $0 = 0 + 0$, sifat distributif
(\cref{thm:g7:priorities:distributivity}) memberi
\[
0 \times a = (0 + 0) \times a = 0 \times a + 0 \times a .
\]
Jadi menambahkan $0 \times a$ pada dirinya sendiri tidak mengubah apa
pun --- dan satu-satunya bilangan yang dapat ditambahkan tanpa
mengubah apa pun adalah $0$. (Secara langsung: kurangkan
$0 \times a$ dari kedua ruasnya.) Karena itu $0 \times a = 0$.

\textbf{2.} Dengan menjabarkan, lalu memakai $1 + (-1) = 0$ dan pertanyaan 1:
\[
a + (-1) \times a
= 1 \times a + (-1) \times a
= \bigl(1 + (-1)\bigr) \times a
= 0 \times a = 0 .
\]
Jadi $(-1) \times a$ adalah bilangan yang, kalau ditambahkan ke $a$,
memberi $0$: itu persis \omterm{def:g7:negatives:def}{lawan} dari $a$, dan $(-1) \times a = -a$.
\Cref{exo:g8:negprod:11} adalah keadaan khusus $a = -1$: di sana,
$(-1) \times (-1) = -(-1) = +1$.

\textbf{3.} Tulislah $-a = (-1) \times a$ (pertanyaan 2) lalu
kelompokkan ulang faktornya:
\[
(-a) \times b
= \bigl((-1) \times a\bigr) \times b
= (-1) \times (a \times b)
= -(a \times b) .
\]

\textbf{4.} Pakailah pertanyaan 3 dua kali (sekali pada tiap faktor):
\[
(-a) \times (-b)
= -\bigl(a \times (-b)\bigr)
= -\bigl(-(a \times b)\bigr)
= a \times b ,
\]
karena \omterm{def:g7:negatives:def}{lawan} dari \omterm{def:g7:negatives:def}{lawan} sebuah bilangan adalah bilangan itu sendiri.

\textbf{5.} Tulislah $a = +d$ atau $a = -d$ dan $b = +e$ atau
$b = -e$, dengan $d$ dan $e$ adalah jaraknya ke \omterm{def:g1:counting:zero}{nol}. Keempat keadaannya:
\[
(+d) \times (+e) = de, \quad
(+d) \times (-e) = -de, \quad
(-d) \times (+e) = -de, \quad
(-d) \times (-e) = +de,
\]
dua yang di tengah menurut pertanyaan 3, yang terakhir menurut
pertanyaan 4. Pada setiap keadaannya, jarak hasilnya adalah
$d \times e$, yaitu \omterm{def:g2:mult:def}{hasil kali} kedua jaraknya, dan tandanya adalah
tanda yang diumumkan tabel pada
\cref{thm:g8:negprod:rules}: tanda sama memberi positif, tanda
berlawanan memberi negatif.

\textbf{6.} $(-1)^2 = +1$, $(-1)^3 = -1$, $(-1)^4 = +1$,
$(-1)^5 = -1$. Secara umum $(-1)^n$ adalah \omterm{def:g2:mult:def}{hasil kali} $n$ faktor
negatif yang jaraknya masing-masing $1$: menurut
\cref{prop:g8:negprod:many} nilainya
$+1$ ketika $n$ \omterm{def:g3:numbers:evenodd}{genap} dan $-1$ ketika $n$ \omterm{def:g3:numbers:evenodd}{ganjil}.

\textbf{7.} Ada sepuluh faktor negatif --- suatu \omterm{def:g3:numbers:evenodd}{bilangan genap} ---
jadi \omterm{def:g2:mult:def}{hasil kalinya} \emph{positif} (\cref{prop:g8:negprod:many}).
Jaraknya adalah $1 \times 2 \times 3 \times \dots \times 10$.

\textbf{8.} Menurut pertanyaan 4, $(-a)^2 = (-a) \times (-a) = a \times a =
a^2$. Lalu menurut pertanyaan 3,
\[
(-a)^3 = (-a)^2 \times (-a) = a^2 \times (-a) = -(a^2 \times a)
= -a^3 .
\]
Tanda minusnya lenyap untuk pangkat \omterm{def:g3:numbers:evenodd}{genap} dan bertahan untuk pangkat
yang \omterm{def:g3:numbers:evenodd}{ganjil}, sejalan dengan \cref{ex:g8:negprod:powers}.

\textbf{9.} Kalau $a = 0$ maka $a^2 = 0$. Kalau tidak, $a$ dan $a$
punya tanda yang \emph{sama}, jadi $a^2 = a \times a$ bernilai positif
(\cref{thm:g8:negprod:rules}). Karena itu kuadrat tidak pernah
negatif; karena $-4$ negatif, tidak ada \omterm{def:g7:negatives:def}{bilangan bertanda} $x$ yang
memenuhi $x^2 = -4$.

\textbf{10.} Zahra benar. Pangkat $2026$ dan $2027$ punya
keganjilan yang berlawanan, jadi $(-1)^{2026} = +1$ dan
$(-1)^{2027} = -1$, yang berjumlah $0$. Secara umum, dari dua pangkat
yang berurutan satu \omterm{def:g3:numbers:evenodd}{genap} dan satu \omterm{def:g3:numbers:evenodd}{ganjil}, jadi kedua nilainya $+1$
dan $-1$ dalam urutan tertentu: bilangan yang saling berlawanan, yang
\omterm{def:g1:addition:def}{jumlahnya} selalu $0$.

\textbf{11.} $A_3 = 1 - 2 + 3 = 2$, lalu
$A_4 = A_3 - 4 = -2$, $A_5 = A_4 + 5 = 3$, $A_6 = A_5 - 6 = -3$,
$A_7 = A_6 + 7 = 4$. Dugaan: $A_n = -\frac{n}{2}$ untuk $n$ \omterm{def:g3:numbers:evenodd}{genap},
dan $A_n = \frac{n+1}{2}$ untuk $n$ \omterm{def:g3:numbers:evenodd}{ganjil}.

\textbf{12.} Untuk $n$ \omterm{def:g3:numbers:evenodd}{genap} sukunya berpasangan seluruhnya:
\[
A_n = (1 - 2) + (3 - 4) + \dots + \bigl((n-1) - n\bigr),
\]
dan tiap pasangannya bernilai $-1$. Ada $\frac{n}{2}$ pasangan (dua
suku pada tiap pasangan, dan $n$ suku seluruhnya), jadi
$A_n = \frac{n}{2} \times (-1) = -\frac{n}{2}$.

\textbf{13.} Untuk $n$ \omterm{def:g3:numbers:evenodd}{ganjil}, \omterm{def:g1:addition:def}{jumlah} $A_n$ adalah \omterm{def:g1:addition:def}{jumlah} $A_{n-1}$
dari $n - 1$ suku pertamanya, ditambah suku terakhirnya, yaitu $+n$
karena $n$ \omterm{def:g3:numbers:evenodd}{ganjil}. Karena $n - 1$ \omterm{def:g3:numbers:evenodd}{genap}, pertanyaan 12 memberi
\[
A_n = A_{n-1} + n = -\frac{n-1}{2} + n
= \frac{-(n-1) + 2n}{2} = \frac{n+1}{2} .
\]

\textbf{14.} $A_{100} = -\frac{100}{2} = -50$, \quad
$A_{101} = \frac{102}{2} = 51$, \quad
$A_{2026} = -\frac{2026}{2} = -1013$.

\textbf{15.} Tidak mengejutkan: berpindah dari $A_{100}$ ke
$A_{101}$ menambahkan tepat satu suku, yaitu $+101$, jadi \omterm{ex:g1:subtraction:difference}{selisihnya}
memang harus $101$. Dengan alasan yang sama,
$A_{2027} - A_{2026} = +2027$ --- tanpa menghitung $A_{2027}$ sama sekali.

\end{solution}
